% Options for packages loaded elsewhere
\PassOptionsToPackage{unicode}{hyperref}
\PassOptionsToPackage{hyphens}{url}
\PassOptionsToPackage{dvipsnames,svgnames,x11names}{xcolor}
%
\documentclass[
  10pt,
  letterpaper,
]{article}
\usepackage{amsmath,amssymb}
\usepackage{iftex}
\ifPDFTeX
  \usepackage[T1]{fontenc}
  \usepackage[utf8]{inputenc}
  \usepackage{textcomp} % provide euro and other symbols
\else % if luatex or xetex
  \usepackage{unicode-math} % this also loads fontspec
  \defaultfontfeatures{Scale=MatchLowercase}
  \defaultfontfeatures[\rmfamily]{Ligatures=TeX,Scale=1}
\fi
\usepackage{lmodern}
\ifPDFTeX\else
  % xetex/luatex font selection
\fi
% Use upquote if available, for straight quotes in verbatim environments
\IfFileExists{upquote.sty}{\usepackage{upquote}}{}
\IfFileExists{microtype.sty}{% use microtype if available
  \usepackage[]{microtype}
  \UseMicrotypeSet[protrusion]{basicmath} % disable protrusion for tt fonts
}{}
\makeatletter
\@ifundefined{KOMAClassName}{% if non-KOMA class
  \IfFileExists{parskip.sty}{%
    \usepackage{parskip}
  }{% else
    \setlength{\parindent}{0pt}
    \setlength{\parskip}{6pt plus 2pt minus 1pt}}
}{% if KOMA class
  \KOMAoptions{parskip=half}}
\makeatother
\usepackage{xcolor}
\usepackage[margin=0.78in]{geometry}
\usepackage{color}
\usepackage{fancyvrb}
\newcommand{\VerbBar}{|}
\newcommand{\VERB}{\Verb[commandchars=\\\{\}]}
\DefineVerbatimEnvironment{Highlighting}{Verbatim}{commandchars=\\\{\}}
% Add ',fontsize=\small' for more characters per line
\newenvironment{Shaded}{}{}
\newcommand{\AlertTok}[1]{\textcolor[rgb]{1.00,0.00,0.00}{\textbf{#1}}}
\newcommand{\AnnotationTok}[1]{\textcolor[rgb]{0.38,0.63,0.69}{\textbf{\textit{#1}}}}
\newcommand{\AttributeTok}[1]{\textcolor[rgb]{0.49,0.56,0.16}{#1}}
\newcommand{\BaseNTok}[1]{\textcolor[rgb]{0.25,0.63,0.44}{#1}}
\newcommand{\BuiltInTok}[1]{\textcolor[rgb]{0.00,0.50,0.00}{#1}}
\newcommand{\CharTok}[1]{\textcolor[rgb]{0.25,0.44,0.63}{#1}}
\newcommand{\CommentTok}[1]{\textcolor[rgb]{0.38,0.63,0.69}{\textit{#1}}}
\newcommand{\CommentVarTok}[1]{\textcolor[rgb]{0.38,0.63,0.69}{\textbf{\textit{#1}}}}
\newcommand{\ConstantTok}[1]{\textcolor[rgb]{0.53,0.00,0.00}{#1}}
\newcommand{\ControlFlowTok}[1]{\textcolor[rgb]{0.00,0.44,0.13}{\textbf{#1}}}
\newcommand{\DataTypeTok}[1]{\textcolor[rgb]{0.56,0.13,0.00}{#1}}
\newcommand{\DecValTok}[1]{\textcolor[rgb]{0.25,0.63,0.44}{#1}}
\newcommand{\DocumentationTok}[1]{\textcolor[rgb]{0.73,0.13,0.13}{\textit{#1}}}
\newcommand{\ErrorTok}[1]{\textcolor[rgb]{1.00,0.00,0.00}{\textbf{#1}}}
\newcommand{\ExtensionTok}[1]{#1}
\newcommand{\FloatTok}[1]{\textcolor[rgb]{0.25,0.63,0.44}{#1}}
\newcommand{\FunctionTok}[1]{\textcolor[rgb]{0.02,0.16,0.49}{#1}}
\newcommand{\ImportTok}[1]{\textcolor[rgb]{0.00,0.50,0.00}{\textbf{#1}}}
\newcommand{\InformationTok}[1]{\textcolor[rgb]{0.38,0.63,0.69}{\textbf{\textit{#1}}}}
\newcommand{\KeywordTok}[1]{\textcolor[rgb]{0.00,0.44,0.13}{\textbf{#1}}}
\newcommand{\NormalTok}[1]{#1}
\newcommand{\OperatorTok}[1]{\textcolor[rgb]{0.40,0.40,0.40}{#1}}
\newcommand{\OtherTok}[1]{\textcolor[rgb]{0.00,0.44,0.13}{#1}}
\newcommand{\PreprocessorTok}[1]{\textcolor[rgb]{0.74,0.48,0.00}{#1}}
\newcommand{\RegionMarkerTok}[1]{#1}
\newcommand{\SpecialCharTok}[1]{\textcolor[rgb]{0.25,0.44,0.63}{#1}}
\newcommand{\SpecialStringTok}[1]{\textcolor[rgb]{0.73,0.40,0.53}{#1}}
\newcommand{\StringTok}[1]{\textcolor[rgb]{0.25,0.44,0.63}{#1}}
\newcommand{\VariableTok}[1]{\textcolor[rgb]{0.10,0.09,0.49}{#1}}
\newcommand{\VerbatimStringTok}[1]{\textcolor[rgb]{0.25,0.44,0.63}{#1}}
\newcommand{\WarningTok}[1]{\textcolor[rgb]{0.38,0.63,0.69}{\textbf{\textit{#1}}}}
\usepackage{longtable,booktabs,array}
\usepackage{calc} % for calculating minipage widths
% Correct order of tables after \paragraph or \subparagraph
\usepackage{etoolbox}
\makeatletter
\patchcmd\longtable{\par}{\if@noskipsec\mbox{}\fi\par}{}{}
\makeatother
% Allow footnotes in longtable head/foot
\IfFileExists{footnotehyper.sty}{\usepackage{footnotehyper}}{\usepackage{footnote}}
\makesavenoteenv{longtable}
\setlength{\emergencystretch}{3em} % prevent overfull lines
\providecommand{\tightlist}{%
  \setlength{\itemsep}{0pt}\setlength{\parskip}{0pt}}
\setcounter{secnumdepth}{5}
\usepackage{amsmath,amssymb,mathtools,booktabs,longtable,array,microtype}
\usepackage[most]{tcolorbox}
\newcommand{\Z}{\mathbb Z}
\newcommand{\C}{\mathbb C}
\newcommand{\Pone}{\mathbb P^1}
\newcommand{\cO}{\mathcal O}
\newcommand{\Pic}{\operatorname{Pic}}
\newcommand{\NS}{\operatorname{NS}}
\newcommand{\rk}{\operatorname{rk}}
\newcommand{\im}{\operatorname{im}}
\newcommand{\coker}{\operatorname{coker}}
\newcommand{\Aut}{\operatorname{Aut}}
\newcommand{\Spec}{\operatorname{Spec}}
\newcommand{\dP}{\mathrm{dP}}
\newtcolorbox{assessmentbox}{colback=gray!4,colframe=black!55,boxrule=.5pt,arc=1pt,left=7pt,right=7pt,top=6pt,bottom=6pt}
\ifLuaTeX
  \usepackage{selnolig}  % disable illegal ligatures
\fi
\usepackage{bookmark}
\IfFileExists{xurl.sty}{\usepackage{xurl}}{} % add URL line breaks if available
\urlstyle{same}
\hypersetup{
  pdftitle={Full Mathematical Peer Review of the 108-Page S6 Manuscript},
  pdfauthor={Independent technical review},
  colorlinks=true,
  linkcolor={blue},
  filecolor={Maroon},
  citecolor={Blue},
  urlcolor={blue},
  pdfcreator={LaTeX via pandoc}}

\title{Full Mathematical Peer Review of the 108-Page S6 Manuscript}
\usepackage{etoolbox}
\makeatletter
\providecommand{\subtitle}[1]{% add subtitle to \maketitle
  \apptocmd{\@title}{\par {\large #1 \par}}{}{}
}
\makeatother
\subtitle{The (3,4,infinity) modular family of 2-tori completed at three
special points}
\author{Independent technical review}
\date{24 August 2026}

\begin{document}
\maketitle

{
\hypersetup{linkcolor=}
\setcounter{tocdepth}{3}
\tableofcontents
}
\section{Referee determination}\label{referee-determination}

\textbf{Recommendation:} Major revision and resubmission after
independent specialist replication. The manuscript should not be
accepted as an established solution in its present form, but this review
did \textbf{not} identify a fatal mathematical contradiction in the
complete 108-page text.

\textbf{Central audit conclusion:} The full manuscript materially
changes the assessment possible from pages 56--66 alone. Section 10
gives a technically substantive proposed explanation of why the printed
arguments of Campana, Demailly, and Peternell (1998, 2020) do not apply
to the constructed nonnormal fibre. The conductor-supported ambient
cotangent section and the monodromy-coinvariant correction are coherent
on close inspection. Thus the earlier packet-only conclusion that the
manuscript was immediately refuted by CDP20 is withdrawn and superseded
by this report.

\textbf{Confidence boundary:} Exact matrix, lattice, Smith-normal-form,
exterior-power, period, and intersection computations were independently
reproduced. The main unresolved audit risk is geometric rather than
arithmetic: the global collapse and integral specialization at the cusp,
the analytic base-change/duality arguments at nonnormal or nonreduced
fibres, and the manuscript's challenge to published literature require
independent checking by specialists. No finite referee audit can convert
an extraordinary new construction into a theorem merely by failing to
locate an error.

The manuscript claims a compact complex threefold \(X\) and a
holomorphic map

\[
f:X\longrightarrow \Pone
\]

whose general fibres are complex two-tori, whose three special fibres
consist of a reduced irreducible nonnormal normal-crossing surface \(W\)
and two multiple fibres \(3S_1\) and \(4S_2\), and for a particular
choice of logarithmic twisting data satisfies

\[
\pi_1(X)=1,\qquad H^*(X;\Z)\cong H^*(S^6;\Z),\qquad X\cong_{\mathrm{diff}}S^6.
\]

It further claims

\[
a(X)=1,\qquad K_X\cong f^*\cO_{\Pone}(-1)\otimes\cO_X(2S_2),
\]

\[
R^1f_*\cO_X\cong\cO\oplus\cO(-1),\qquad R^2f_*\cO_X\cong\cO(-1),
\]

and \(c_3(X)=2\). These conclusions would resolve the classical Hopf
problem in the affirmative and, at the same time, contradict the printed
conclusions of Campana et al.~(2020, Theorem 2.2 and Corollary 2.3).
That combination makes the burden of proof exceptionally high.

My disposition is therefore deliberately two-level:

\begin{enumerate}
\def\labelenumi{\arabic{enumi}.}
\tightlist
\item
  \textbf{Internal mathematical audit:} I found the proof architecture
  coherent and found no theorem-breaking inconsistency. The exact
  computations tested independently are correct.
\item
  \textbf{Publication-level validation:} The manuscript is not yet ready
  to be treated as a settled proof. Several pivotal, nonstandard
  geometric arguments require expansion and independent specialist
  replication, especially because Section 10 would invalidate published
  theorem statements for the proposed example.
\end{enumerate}

\section{Scope, source status, and review
method}\label{scope-source-status-and-review-method}

\subsection{Source completeness and
provenance}\label{source-completeness-and-provenance}

The reviewed archive, \texttt{s6\_canonical\_108\_pages\_audited.zip},
is complete as a reconstruction package. It contains ten modular source
parts, a modular assembler, a monolithic source, a 108-page compiled
PDF, an authority reassembly, build scripts, checksum manifests,
page-by-page visual audits, and reconstruction provenance. The archive's
own verification script passes. The modular and monolithic sources
compile to byte-identical PDFs in the audited environment.

This establishes fidelity to the supplied authority PDF, not authorship
of the TeX source and not mathematical correctness. The archive
expressly states that it is an independent reconstruction rather than
the author's original TeX. For any formal submission, an
author-controlled source or an explicit author sign-off on the
reconstructed canonical PDF is necessary.

The earlier review package was based on an incomplete 155-line assembler
plus packet 06. Its conclusion that the main theorem was directly
contradicted by CDP20 was a correct description of the logical conflict
visible in that slice, but an incorrect final assessment of the complete
manuscript because Section 10 was unavailable. This report replaces that
review in full.

\subsection{Review procedures}\label{review-procedures}

The audit used four complementary procedures.

First, I read the complete proof dependency chain rather than reviewing
isolated claims. In particular, I traced the construction from the
monodromy representation through the period family, the three local
fillings, the global gluing, both integral topology computations, smooth
recognition, analytic invariants, and the literature-conflict analysis.

Second, I independently reproduced exact computations with integer and
symbolic arithmetic. The supplied audit package contains
\texttt{audit/exact\_checks.py} and a machine-readable result file.
Twenty-eight checks pass. These include monodromy orders, contragredient
matrices, fixed lattices, the invariant alternating form, Smith normal
forms, exterior-power monodromy, Wang ranks, saturation tests, the
fundamental-group determinant, the period transformation identities, the
Neron-Severi polynomial, and the Hermitian Gram determinant.

Third, I compared the two topology routes. The first route uses an
explicit collapse and Mayer--Vietoris. The second uses the integral
Leray spectral sequence and Appendix B's nearby-cycle specialization
theorem. Agreement between genuinely different calculations is
significant evidence, although it is not a substitute for verifying
their shared geometric inputs.

Fourth, I reviewed Section 10 as an argument about the exact printed
proof mechanisms in Campana et al.~(1998, 2020), not merely as an
assertion that those papers are wrong. The review distinguishes: (a) the
formal contradiction between conclusions; (b) the manuscript's proposed
failure mechanism; and (c) whether that proposed mechanism has been
independently validated.

\subsection{Meaning of audit labels}\label{meaning-of-audit-labels}

``Verified computationally'' means that the displayed finite calculation
has been reproduced exactly. ``No contradiction found'' means that the
proof was read and its logical transitions appear coherent, but it is
not a certification. ``Major revision request'' identifies a pivotal
argument whose present exposition is too compressed for an extraordinary
theorem, even where no specific counterexample has been found.
``Confirmed error'' is reserved for a statement or proof sentence for
which a definite defect can be exhibited.

\section{The proof architecture and its dependency
graph}\label{the-proof-architecture-and-its-dependency-graph}

The main theorem has the following dependency chain.

\begin{enumerate}
\def\labelenumi{\arabic{enumi}.}
\tightlist
\item
  Sections 2--3 construct a rank-four integral local system over the
  \((3,4,\infty)\) orbifold and a nonpolarized holomorphic family of
  complex two-tori with monodromies of orders \(3\), \(4\), and
  unipotent type at the cusp.
\item
  Section 4 constructs a toric filling at the cusp from an infinite
  \(A_2\) fan and a free, properly discontinuous \(\Z^2\) action. Its
  central fibre \(W\) is irreducible and nonnormal, with normalization
  \(\dP_6\), three double curves, and two triple points.
\item
  Section 5 constructs the order-\(3\) and order-\(4\) logarithmic
  fillings, with reduced fibres bielliptic surfaces.
\item
  Section 6 glues the smooth family and three fillings into a compact
  complex threefold \(X\to\Pone\) and records the discrete gluing
  parameter \(\ell_0\).
\item
  Section 7 computes \(\pi_1(X)\) and integral homology. For the chosen
  twists \((\ell_0,\ell_1,\ell_2)=(0,1,-1)\), the controlling integer \[
  p=12\ell_0-4\ell_1-3\ell_2
  \] is \(-1\). The manuscript gives two integral homology computations.
\item
  Section 8 uses the simply connected integral homology-sphere
  conclusion, Hurewicz--Whitehead theory, and \(\Theta_6=0\) to identify
  \(X\) diffeomorphically with \(S^6\).
\item
  Section 9 computes algebraic dimension, direct images, the canonical
  bundle, Chern-number consequences, holomorphic forms, and
  automorphisms.
\item
  Section 10 explains why the printed CDP98/CDP20 proof mechanisms are
  claimed not to cover this \(X\).
\end{enumerate}

A fatal defect in Sections 2--6 would invalidate the space. A defect in
either topology route could potentially be repaired by the other only if
their shared local specialization and gluing data remain valid. A defect
in Section 9 would not by itself destroy the complex structure on
\(S^6\), but it could remove the claimed contradiction with CDP20. A
defect in Section 10 would reinstate the literature obstruction unless
the published theorem itself were independently corrected.

\section{Section-by-section mathematical
review}\label{section-by-section-mathematical-review}

\subsection{Lattice and monodromy data (Section
2)}\label{lattice-and-monodromy-data-section-2}

The matrices

\[
T_1=\begin{pmatrix}
1&0&-6&2\\0&-1&1&1\\0&-1&0&1\\0&0&0&1
\end{pmatrix},\qquad
T_2=\begin{pmatrix}
1&6&0&-3\\0&0&-1&1\\0&1&0&0\\0&0&0&1
\end{pmatrix}
\]

satisfy

\[
T_1^3=I,\qquad T_2^4=I,\qquad T_0=(T_1T_2)^{-1}=I+N,\qquad N^2=0.
\]

The exact orders are indeed \(3\) and \(4\). The displayed
contragredients \(A_j=(T_j^{-1})^t\) and \(M_0=(T_0^{-1})^t\) are
correct, and \(A_1A_2M_0=I\). The invariant and vanishing-cycle lattices
are also correct:

\[
\ker(M_0-I)=\im(M_0-I)=\langle\widehat w,\widehat\delta\rangle,
\]

and the induced map

\[
B_0:\overline\Lambda\to\Lambda_{\mathrm{tor}},\qquad
B_0=\begin{pmatrix}0&1\\-1&0\end{pmatrix},
\]

is unimodular. The fixed lattices

\[
\Lambda^{A_1}=\langle\varepsilon,\widehat\delta\rangle,
\qquad
\Lambda^{A_2}=\langle\varepsilon',\widehat\delta\rangle
\]

are primitive rank-two sublattices with the stated generators.

The invariant alternating form \(Q_0\) is also consistent. In the
ordered basis \((\gamma,u,w,\delta)\), its nonzero entries are
\(Q_0(\gamma,\delta)=1\) and \(Q_0(u,w)=6\). Direct multiplication gives
\(T_j^tQ_0T_j=Q_0\). The induced form

\[
b_N(x,y)=Q_0(x,Ny)
\]

on \(V/\ker N\), in the basis \(\bar w,\bar\delta\), is
\(\operatorname{diag}(6,-1)\). This confirms the manuscript's central
nonpolarized feature: the monodromy-compatible alternating form has
indefinite limiting form and does not define a polarization. This
matters because standard polarized degeneration theorems cannot simply
be imported (Clemens, 1977; Mumford, 1972).

The group-theoretic uniformization by the \((3,4,\infty)\) triangle
group is standard. The manuscript carefully separates the uniformizing
coordinate from the period coordinate at the cusp and derives the
translation sign. I found no inconsistency in the presentation
\(\Delta\cong\Z/3*\Z/4\) or the definition of the representation.

\textbf{Disposition:} Exact algebra verified. No correction required,
except that a concise table of all actions on \(V\), \(\Lambda\), and
their exterior powers would reduce later sign risk.

\subsection{Period family (Section 3)}\label{period-family-section-3}

The period matrix is built from

\[
\sigma_1=6\mu\gamma+\tau u+w,
\qquad
\sigma_2=\beta\gamma+\mu u+\delta,
\]

with \(F^1=\langle\sigma_1,\sigma_2\rangle\). The construction proceeds
through three nontrivial analytic steps: a modular function \(\tau\), a
torsor solution \(\mu\) for an automorphy sheaf of degree \(-1\), and a
solution \(\beta\) to a further additive cocycle. The final imaginary
constant is chosen so that the lattice remains nondegenerate.

The modular transformation laws used later are mutually consistent. The
automorphy sheaf calculation in Lemma 3.10, including the degree \(-1\)
conclusion, is coherent: the chosen modular expression has the required
character, and its divisor behavior at the cusp gives
\(\mathscr L\simeq\cO(-1)\). Because a square-root factor is used, the
argument would be clearer if formulated invariantly in terms of the
orbifold line bundle and its character, rather than by an implicitly
chosen branch. I did not find a contradiction in the present
construction.

The nondegeneracy argument is important. It is not enough that
\(\Pi(z)\) be formally equivariant; its four real columns must span a
full lattice in \(\C^2\). The manuscript reduces the determinant to an
inequality involving

\[
D=\operatorname{Im}\beta-\frac{6(\operatorname{Im}\mu)^2}{\operatorname{Im}\tau}
\]

and uses the free imaginary constant in \(\beta\) to force \(D<0\)
throughout a compact fundamental domain. The proper-discontinuity and
compactness input makes this strategy plausible. The same expression
later controls the signature of the Hermitian form associated with
\(\eta\).

The cusp asymptotic decomposition

\[
Z=sB_0+C(t_c)
\]

is the precise bridge to the toric filling. I checked the equivariance
and the exponential identity that produces the torus element
\(c_{\bar\lambda}(t_c)\). No algebraic inconsistency was found.

\textbf{Disposition:} No fatal defect found. For publication, the torsor
existence and the global nondegeneracy choice should be stated as
standalone propositions with explicit cohomology groups and norm
estimates. The current proof is credible but densely packed.

\subsection{Toric filling at the cusp (Section
4)}\label{toric-filling-at-the-cusp-section-4}

The fan consists of cones over the cells of the periodic \(A_2\)
triangulation at height one. Every maximal cone is unimodular because
the edge differences form a \(\Z^2\) basis. Local finiteness of the
triangulation gives a Hausdorff, second-countable toric complex
manifold, and the character \(t=\chi^{(0,0,1)}\) has reduced
normal-crossing divisor

\[
t^{-1}(0)=\bigcup_{v\in\Z^2}E_v.
\]

The action

\[
\Psi_{\bar\lambda}=\tau_{\bar\lambda}\circ\Phi_{\bar\lambda}
\]

combines integral translation of the fan with a \(t\)-dependent torus
translation. The manuscript does more than assert properness: it
introduces a rescaled position map and obtains a quantitative
displacement estimate. Since \(B_0\) is orthogonal and unimodular, large
group elements move compact sets apart uniformly near the cusp. The
compact fundamental-set argument then yields properness of the quotient
map to the disc.

The action is free. On the central fibre, the fan translation is free on
irreducible components; away from the central fibre, the action
corresponds to a free lattice translation on the torus cover. The
quotient therefore has local models

\[
t=z_0,
\qquad t=z_0z_1,
\qquad t=z_0z_1z_2,
\]

and is smooth. The invariant logarithmic/top form gives the asserted
relative canonical trivialization.

The normalization and conductor combinatorics are consistent. A
fundamental component is \(\dP_6\); opposite sides of its toric hexagon
are identified. This yields three rational double curves and two triple
points. The orbit stratification

\[
W=(\C^*)^2\sqcup 3\C^*\sqcup\{P,Q\}
\]

gives \(e(W)=2\).

I found one definite exposition defect. Proposition 4.7(iv) says the
image of the section is closed because it is ``proper over \(\Delta\).''
That is not the correct reason when the base is an open disc. The
conclusion is readily repaired: a continuous section of a Hausdorff map
has closed graph/image, or one may prove closedness directly in the
chosen quotient chart. This does not threaten the construction.

The highest-risk part of Section 4 is not a visible algebraic error but
the global topological bookkeeping for an infinite fan and quotient. The
manuscript contains the needed ideas, yet a publication version should
consolidate them into one theorem: local finiteness, second
countability, freeness, proper discontinuity, compact fundamental
domain, and compatibility of the toric deformation with the group
action.

\textbf{Disposition:} No theorem-breaking defect found. One minor proof
sentence should be corrected; global quotient topology deserves
expansion.

\subsection{Finite-order logarithmic fillings (Section
5)}\label{finite-order-logarithmic-fillings-section-5}

The local family over the linearizing disc carries an affine
transformation whose linear part is \(A_j\) and whose translation part
is \(v_j/m_j\). The admissibility conditions ensure that the affine
cyclic action is free. The fixed-point criterion is correctly reduced to
membership of the translation vector in \((A_j-I)\) of the torus
lattice, and the chosen vectors avoid that image.

The quotient produces a smooth threefold with multiple fibre \(m_jS_j\).
The reductions \(S_j\) are bielliptic surfaces. The computed first
homology and torsion are compatible with the standard Bagnera--de
Franchis classification (Barth et al., 2004). The order-\(3\) and
order-\(4\) cases also generate the precise local specialization indices
later used in Section 7.

The logarithmic regluing section is arranged so that the transformed
zero section is invariant. The conjugation identity used to compare the
punctured family with the quotient is correct. The resulting drift
controls \(\ell_j=\gamma(v_j)\).

\textbf{Disposition:} No fatal error found. The proof would be more
auditable with a single diagram showing the cover, affine action,
quotient, multiple fibre, reduction, and boundary identification. The
classification of the two bielliptic surfaces should be tied explicitly
to the standard notation in Barth et al.~(2004).

\subsection{Global gluing and compactness (Section
6)}\label{global-gluing-and-compactness-section-6}

The three local models are glued to the smooth family over \(B^\circ\).
The overlap maps are biholomorphisms by construction. The manuscript
addresses Hausdorffness through shrinkable collars and the fact that
each local piece maps properly to a disc. Properness of the global map
is local on the base for a separated holomorphic map; the resulting
total space is compact because the base is compact and the map is
proper.

The zero section is used as more than notation: it fixes the lift of
boundary meridians and prevents hidden fibre translations in the
finite-order gluings. At the cusp, affine regluing changes the discrete
parameter \(\ell_0\), and the tautological gluing gives \(\ell_0=0\).

No incompatibility was found among the boundary identifications. The
proof is nevertheless unusually sensitive to the distinction between the
smooth fibre lattice, the toric translation lattice, and the quotient of
\(\Lambda\) by \(\Lambda_{\mathrm{tor}}\). A one-page gluing table
should identify, for every boundary component, the deck group, meridian,
fibre translation, orientation, and induced relation in \(\pi_1\).

\textbf{Disposition:} Coherent. Major revision only in presentation and
auditability, not because of a demonstrated flaw.

\subsection{Fundamental group (Section
7.5)}\label{fundamental-group-section-7.5}

After attaching the cusp filling, the normal closure of
\(\Lambda_{\mathrm{tor}}\) is \(\ker\gamma\). Hence the image of the
fibre group is generated by one central element \(c\). The two finite
fillings impose

\[
x^3=c^{\ell_1},\qquad y^4=c^{\ell_2},
\]

and the cusp relation gives \(xy=c^{\ell_0}\). Thus

\[
\pi_1(X)=\langle c,x,y\mid c\text{ central},\ xy=c^{\ell_0},\ x^3=c^{\ell_1},\ y^4=c^{\ell_2}\rangle.
\]

Eliminating \(y\) makes the group abelian with relation matrix

\[
R=\begin{pmatrix}
3&-\ell_1\\
4&\ell_2-4\ell_0
\end{pmatrix},
\qquad
\det R=-(12\ell_0-4\ell_1-3\ell_2).
\]

Because the entries include \(3\) and \(4\), the first Smith invariant
is \(1\); hence

\[
\pi_1(X)\cong\Z/|p|,
\qquad
p=12\ell_0-4\ell_1-3\ell_2.
\]

The admissibility congruences imply \(\gcd(p,12)=1\). For \((0,1,-1)\),
\(p=-1\) and the group is trivial. For the comparison sign, \(p=-7\).

These calculations pass exact Smith-normal-form checks. The manuscript's
long sign lemma is not redundant: a relative sign error would
interchange the simply connected and \(\Z/7\) examples. I found its
treatment internally consistent. The final invariant \(|p|\) is
orientation-independent, but the manuscript should provide a compact
convention table because the proof currently requires the reader to
track clockwise meridians through several covers.

\textbf{Disposition:} Algebra verified. No mathematical correction
required.

\subsection{First integral homology route: collapse and Mayer--Vietoris
(Sections
7.1--7.6)}\label{first-integral-homology-route-collapse-and-mayervietoris-sections-7.17.6}

The first route constructs an equivariant collapse of a nearby torus
fibre onto the singular fibre. The central fibre has homology

\[
H_k(W;\Z)=(\Z,\Z^2,\Z^4,\Z^2,\Z),
\qquad k=0,\ldots,4.
\]

The boundary mapping torus has Wang ranks

\[
(1,3,6,6,3,1).
\]

The kernel of its map to \(W\) is described by the toric meridian and
the two vanishing-cycle tori. Attaching the remainder of the fibration
and computing the relevant quotient maps gives, for admissible data,

\[
H_k(X;\Z)\cong\Z/|p|\quad (1\le k\le4),
\]

with \(H_0=H_6=\Z\) and the remaining groups determined by duality. For
\(p=\pm1\), the middle homology vanishes.

The exterior-power monodromy and all displayed lattice ranks are
correct. In particular, for

\[
M_0=\begin{pmatrix}
1&0&0&0\\0&1&0&0\\0&1&1&0\\-1&0&0&1
\end{pmatrix},
\]

the Wang ranks are exactly \((1,3,6,6,3,1)\). The image of
\(\bigwedge^2M_0-I\) is the saturated rank-two lattice generated by the
stated vectors. The matrix called \(\alpha_2\) has rank \(3\), gcd of
its nonzero \(3\times3\) minors equal to \(1\), and the displayed
primitive left-kernel generator.

The unresolved risk is geometric: Proposition 7.2's global collapse is
elaborate. The manuscript constructs it using polar coordinates, a
honeycomb subdivision, and lifted deformation retractions. It appears
designed to address continuity across double and triple strata and
equivariance under the toric lattice. I did not locate a discontinuity
or an equivariance failure. Nevertheless, this is one of the places
where an independent topology specialist should reconstruct the map
rather than merely inspect the prose. A cell-level chain model or an
explicit diagram would substantially improve verifiability.

\textbf{Disposition:} All finite algebra checks pass. The collapse
theorem remains a major specialist-verification item.

\subsection{Second integral homology route: Leray spectral sequence
(Section
7.7)}\label{second-integral-homology-route-leray-spectral-sequence-section-7.7}

The second route is structurally independent. Let
\(j:B^\circ\hookrightarrow B\) and let \(\mathcal L^q\) be the local
system with fibre \(\bigwedge^qV\). The manuscript identifies
\(R^qf_*\Z\) as a subsheaf of \(j_*\mathcal L^q\), with local quotient
skyscrapers determined by specialization at the three special fibres.

At the finite fibres the pullback lattices have the claimed indices. In
degree one, the indices are \(3\) and \(4\). At the order-\(4\) fibre
there is an additional index \(2\) in degrees two and three. The
degree-two conclusion follows from the invariant-lattice Gram
determinant and evenness of the bielliptic intersection form. The
degree-three index follows by pairing with the degree-one image. These
calculations are correct at the lattice level.

At the cusp, Theorem B.1 asserts that specialization is injective and
its image is the full \(T_0\)-invariant lattice in every degree.
Consequently the global sections are

\[
H^0(B,R^1)=12\Z\gamma,
\quad
H^0(B,R^2)=\Z\cdot2q,
\quad
H^0(B,R^3)=\Z\cdot2\gamma uw,
\quad
H^0(B,R^4)=12\Z\operatorname{vol}.
\]

The coinvariant groups in degree two are torsion-free, with
\([uw]=6[\gamma\delta]\). The parabolic cohomology groups \(H^1(B,R^1)\)
and \(H^1(B,R^2)\) vanish by the displayed intersection calculations.

Multiplicativity reduces all nonzero \(d_2\) maps to the same integer,
up to sign. The geometric extension class is
\(p=12\ell_0-4\ell_1-3\ell_2\). Hence, for \(p=\pm1\), the spectral
sequence yields the cohomology of \(S^6\).

I independently checked the arithmetic underlying the multiplicative
argument. The difficult input is not the final differential calculation
but the integral identification of the cusp specialization lattice. That
input is reviewed separately below.

\textbf{Disposition:} Internally consistent and strongly corroborative.
Conditional on Appendix B, the second route is persuasive.

\subsection{Recognition as the smooth six-sphere (Section
8)}\label{recognition-as-the-smooth-six-sphere-section-8}

Once \(X\) is simply connected and has the integral homology of \(S^6\),
the Hurewicz theorem and Whitehead's theorem make it a homotopy
\(6\)-sphere. Smooth homotopy \(6\)-spheres are classified by
\(\Theta_6\), and \(\Theta_6=0\) (Kervaire \& Milnor, 1963). Thus the
diffeomorphism conclusion follows.

This step is standard and correct, conditional on Section 7. The
manuscript could make explicit that the smooth structure is oriented by
the complex orientation, although in dimension six there is no exotic
sphere ambiguity.

\textbf{Disposition:} Valid conditional deduction.

\subsection{Algebraic dimension (Section
9.1)}\label{algebraic-dimension-section-9.1}

The Neron--Severi calculation begins with the condition

\[
E\wedge\sigma_1\wedge\sigma_2=0.
\]

For a general integral two-form \(E\), the manuscript obtains a
polynomial \(P_E\) in \(1,\tau,\mu,\beta\), and \(6\mu^2-\tau\beta\).
The symbolic expansion is correct. Lemma 9.3 proves these five
holomorphic functions are linearly independent by applying the \(g_0\)
and \(g_2\) transformation laws. I reproduced the successive-difference
identities.

It follows that the very general fibre has

\[
\NS(F_b)=\Z\eta,
\qquad
\eta=u\wedge w+6\gamma\wedge\delta.
\]

The square is

\[
\eta^2=12.
\]

The associated Hermitian Gram matrix is

\[
G_\eta=
\begin{pmatrix}
2\operatorname{Im}\tau&2\operatorname{Im}\mu\\
2\operatorname{Im}\mu&\frac13\operatorname{Im}\beta
\end{pmatrix},
\]

with determinant

\[
\det G_\eta=
\frac23\operatorname{Im}\tau
\left(
\operatorname{Im}\beta-
\frac{6(\operatorname{Im}\mu)^2}{\operatorname{Im}\tau}
\right)<0.
\]

Thus \(\eta\) has signature \((1,1)\) and no nonzero multiple is
semipositive. The manuscript concludes that the very general torus has
algebraic dimension zero. This is consistent with the standard theory of
divisors on complex tori: effective divisors induce semipositive
Hermitian forms, while algebraic dimension one for a two-torus is
governed by a nonzero isotropic semipositive class (Birkenhake \& Lange,
2004; Lange \& Birkenhake, 1999). Because the proof compresses this
classification step, it should state the exact proposition being used.

A meromorphic function on \(X\) restricts to a constant on very general
fibres. The manuscript then proves that constancy on uncountably many
fibres forces descent to the base. Hence

\[
\mathcal M(X)=f^*\mathcal M(\Pone),
\qquad a(X)=1,
\]

and \(f\) is the algebraic reduction.

\textbf{Disposition:} No contradiction found; exact period algebra
verified. Add an explicit torus-classification citation or proof.

\subsection{Canonical bundle and direct images (Sections
9.2--9.4)}\label{canonical-bundle-and-direct-images-sections-9.29.4}

The relative two-form \(e=d\zeta_1\wedge d\zeta_2\) transforms by
\(\det R_g\). Its behavior at the finite orbifold points gives
fractional vanishing orders \(2/3\) and \(1/4\), and it is regular and
nonvanishing at the cusp. The resulting direct image is

\[
f_*\omega_{X/B}\cong\cO_{\Pone}(1).
\]

The multiple-fibre correction yields

\[
K_X\cong f^*\cO(-1)\otimes\cO(2S_2).
\]

The character arithmetic is coherent. The non-torsion conclusion follows
because the stated bundle has a nontrivial multiple-fibre component and
its appropriate power is a nontrivial pullback.

For the structure-sheaf direct images, the manuscript obtains

\[
f_*\cO_X=\cO,
\qquad
R^1f_*\cO_X\cong\cO\oplus\cO(-1),
\qquad
R^2f_*\cO_X\cong\cO(-1).
\]

The fibre cohomology at \(W\) is computed by the normalization/conductor
sequence. The finite multiple fibres are handled by equivariant local
models and descending induction through the nilpotent thickening.
Relative duality identifies the top direct image. The eigenweights
determine the splitting of \(R^1\).

I found no algebraic inconsistency, and the resulting global Hodge
numbers agree with the Leray calculation. The main concern is the
precision of the analytic foundations. In a theorem of this importance,
references to Grauert, base change, and relative duality should be
attached to exact statements and hypotheses. The reader should not have
to infer that the scheme-theoretic multiple fibres are covered, that the
relevant sheaves are flat, and that duality applies to the Gorenstein
nonnormal surface \(W\).

\textbf{Disposition:} Plausible and internally consistent. Major
revision requested for exact analytic theorem statements.

\subsection{Holomorphic forms, automorphisms, and Chern data (Sections
9.5--9.7)}\label{holomorphic-forms-automorphisms-and-chern-data-sections-9.59.7}

The vanishing of global holomorphic \(p\)-forms is derived by writing
invariant forms on the universal cover and using the period derivatives
and local growth at the cusp and orbifold points. The weight-zero and
weight-one lemmas eliminate the remaining coefficients. The argument is
technically detailed and consistent with the computed \(h^{0,q}\).

The identity component of the automorphism group must preserve the
algebraic reduction and hence acts fibrewise. Translation vector fields
correspond to the invariant lattice line
\(\Lambda^\Delta=\Z\widehat\delta\), giving a one-dimensional vertical
field. Its flow is \(\C^*\) rather than an elliptic translation group
because of the degeneration. The toric fixed set is one rational double
curve with normal weights \((+1,-1)\). This is compatible with
\(e(X^{\C^*})=e(X)=2\).

Because \(H^2(X;\Z)=H^4(X;\Z)=0\), \(c_1\) and \(c_2\) vanish
integrally. The top Chern number equals the Euler characteristic,

\[
\langle c_3(T_X),[X]\rangle=e(X)=2.
\]

The Riemann--Roch arithmetic giving \(\chi(\cO_X)=0\),
\(\chi(\Omega_X^1)=-1\), and \(\chi(T_X)=1\) is consistent.

\textbf{Disposition:} No contradiction found. These conclusions remain
conditional on the topology and direct-image calculations.

\section{Appendix B: integral specialization at the nonnormal
cusp}\label{appendix-b-integral-specialization-at-the-nonnormal-cusp}

Appendix B is one of the two most important parts of the manuscript. It
is intended to remove reliance on the geometric collapse for the
statement

\[
\im\bigl(H^q(W;\Z)\to H^q(F;\Z)\bigr)
=H^q(F;\Z)^{T_0}
\]

and to establish injectivity and saturation integrally.

\subsection{Nearby-cycle sheaves}\label{nearby-cycle-sheaves}

The local normal-crossing models give Milnor fibres of homotopy type
point, \(S^1\), and \(T^2\). Therefore

\[
\mathcal H^0\psi=\Z_W,
\qquad
\mathcal H^1\psi=\mathcal V,
\qquad
\mathcal H^2\psi=\Z_P\oplus\Z_Q.
\]

The global orientation convention along the three double curves
determines the restriction maps at the two triple points. In both cases
the image consists of triples satisfying

\[
n_1-n_2+n_3=0.
\]

Thus the exact sequence

\[
0\to\mathcal V\to\bigoplus_{k=1}^3\iota_{k*}\Z_{\bar C_k}
\xrightarrow{\delta_1}\Z_P\oplus\Z_Q\to0
\]

is correct provided the branch labels and orientations are those fixed
in the normalization. The manuscript performs the sign check explicitly
at both triple points.

\subsection{\texorpdfstring{The \(E_2\)
page}{The E\_2 page}}\label{the-e_2-page}

The normalization/conductor resolution gives

\[
H^\bullet(W;\Z)=(\Z,\Z^2,\Z^4,\Z^2,\Z),
\]

all free. The calculation of \(H^\bullet(W;\mathcal V)\) gives

\[
(\Z^2,\Z,\Z^3,0),
\]

and the top nearby-cycle row contributes \(\Z^2\) in bidegree \((0,2)\).
The total rank on the \(E_2\) page is \(18\), while the smooth torus
fibre has total Betti number \(16\). Hence exactly one rank must be
killed.

The degree-by-degree rank bookkeeping forces all differentials touching
the bottom row to vanish except for the one needed in total degree
three. This part is internally consistent. Since all groups are free,
rank information does not by itself prove saturation; the appendix
addresses that separately.

\subsection{Saturation and the degree-three Gysin
argument}\label{saturation-and-the-degree-three-gysin-argument}

The key additional input is a compactly supported cycle transported
along a torus orbit. The manuscript constructs a local trivialization
near that orbit, uses the Gysin map to identify a primitive degree-three
class, and shows that specialization has rank two and saturated image in
degree three. Poincare duality and multiplicativity then propagate the
saturation and invariant-lattice identification to all degrees.

I did not locate a rank mismatch, sign inconsistency, or circular
dependency. Appendix B explicitly states which parts still rely on the
geometric collapse: the identification of particular generators and
normalizations, not the injectivity and invariant-image theorem used by
the Leray route.

The revision problem is exposition and independent reproducibility. The
argument should state the nearby-cycle functor, the specialization
morphism, the edge map, and the compact-support/Gysin construction in a
standard derived-category framework, with precise integral references
such as Deligne's SGA 7 formalism or a modern equivalent (Deligne, 1973;
Dimca, 2004; Kashiwara \& Schapira, 1990). The present argument is
plausible, but it is too consequential to leave any ambiguity about
conventions or functoriality.

\textbf{Disposition:} No fatal defect found; highest-priority specialist
verification item.

\section{Section 10 and the conflict with
CDP98/CDP20}\label{section-10-and-the-conflict-with-cdp98cdp20}

\subsection{The formal contradiction}\label{the-formal-contradiction}

If the manuscript's main conclusions hold, then \(X\) satisfies

\[
H^1(X;\Z)=H^2(X;\Z)=0,
\qquad a(X)=1,
\]

and has a holomorphic algebraic reduction \(f:X\to\Pone\). Campana et
al.~(2020, Theorem 2.2) print, under those hypotheses, \(c_3(X)\le0\),
and their Corollary 2.3 prints \(a(X)=0\) for a compact complex
threefold homeomorphic to \(S^6\). The manuscript instead has
\(c_3(X)=2\) and \(a(X)=1\). Both sets of conclusions cannot
simultaneously be correct.

The complete manuscript does not evade this issue. It isolates two exact
proof mechanisms and argues that one is irreparably false at the
nonnormal fibre \(W\), while the other is a monodromy error that can be
repaired for this example.

\subsection{The nonnormal-fibre cotangent
mechanism}\label{the-nonnormal-fibre-cotangent-mechanism}

Let \(S\) be a reduced normal-crossing fibre in a smooth threefold
\(\mathcal X\to\Delta\), with normalization \(\eta:\widetilde S\to S\)
and conductor \(D\). Locally the defining function is \(g=ux\), \(uxy\),
or \(uxyz\). The differential \(dg|_S\) is nonzero away from the
conductor but restricts to zero on the conductor.

For a line bundle \(A\) on \(S\) whose pullback to \(\widetilde S\) is
trivial, choose a nowhere-zero pullback section \(e\). The
normalization/conductor Mayer--Vietoris sequence lets the pair

\[
\bigl(\eta^*(dg|_S)\otimes e,0\bigr)
\]

descend to a nonzero section

\[
s\in H^0(S,\Omega^1_{\mathcal X}|_S\otimes A).
\]

Its image in \(\Omega_S^1\otimes A\) is supported on the conductor and
is torsion. Therefore its image in the torsion-free quotient

\[
\widetilde\Omega_S^1=\Omega_S^1/\mathcal T(\Omega_S^1)
\]

is zero. If \(A\) is nontrivial, \(s\) does not arise from the conormal
line. This produces exactly the type of ambient cotangent section that
is invisible to a vanishing statement for
\(\widetilde\Omega_S^1\otimes A\).

This is the strongest part of the manuscript's rebuttal. The local
formulas are elementary and the descent through the conductor sequence
is explicit. The mechanism shows that the implication

\[
H^0(S,\widetilde\Omega_S^1\otimes A)=0
\quad\Longrightarrow\quad
H^0(S,\Omega^1_X|_S\otimes A)=0
\]

is not valid for nonnormal \(S\) merely from vanishing of the conormal
contribution. For normal \(S\), extension across codimension two can
restore the argument; for the nonnormal surface \(W\), the conductor
supports the obstruction.

For the constructed \(W\), every line bundle from \(X\) pulls back
trivially to the normalization \(\dP_6\) because \(c_1=0\) and
\(\Pic(\dP_6)\hookrightarrow H^2(\dP_6;\Z)\). Serre--Grothendieck
duality then gives

\[
H^2\bigl(W,(T_X\otimes L)|_W\bigr)^*
\cong
H^0\bigl(W,\Omega_X^1|_W\otimes(L^*\otimes\omega_X)|_W\bigr)
\ne0.
\]

Top-degree cohomology and base change imply
\(R^2f_*(T_X\otimes L)\ne0\). The manuscript therefore claims that the
hypothesis used by Campana et al.~(2020, Proposition 2.4) holds for no
\(L\) on this \(X\).

I found no local algebra error in this mechanism. The remaining burden
is to state the exact analytic duality and top-degree base-change
theorems with all hypotheses. Because this argument would invalidate a
published theorem's application, it must be independently refereed by
experts in nonnormal complex spaces and analytic duality.

\subsection{The monodromy-coinvariant
correction}\label{the-monodromy-coinvariant-correction}

The second issue concerns a split group extension

\[
1\to K\to G\to Q\to1.
\]

The correct abelianization is

\[
G^{\mathrm{ab}}\cong Q^{\mathrm{ab}}\oplus(K^{\mathrm{ab}})_Q,
\]

where the second summand consists of coinvariants under conjugation. It
is not generally \(Q^{\mathrm{ab}}\oplus K^{\mathrm{ab}}\) unless the
action is trivial. This is standard group homology.

For the manuscript's torus fibration, the unipotent monodromy acts
nontrivially on \(H_1(F;\Z)\). The coinvariant rank is one, not four.
Consequently, formulas in the printed CDP98/CDP20 route that use
\(b_1(F)\) rather than the coinvariant rank do not apply to this
example. This also explains how

\[
H^1(X,\cO_X)\cong\C
\]

can coexist with

\[
H^1(F,\cO_F)\cong\C^2
\]

without surjectivity of restriction.

The manuscript then proves by a separate countability argument that a
line bundle non-torsion on every fibre nevertheless exists. Thus the
monodromy problem breaks the printed route to that input but can be
repaired for this example. By contrast, the nonnormal cotangent
reduction remains unavailable.

The group-theoretic correction is correct. Its application to the exact
sequence in CDP98 should still be checked against the precise local
topological model and notation in the original paper, but the
mathematical principle is not in doubt.

\subsection{Referee assessment of the literature
dispute}\label{referee-assessment-of-the-literature-dispute}

The earlier packet-only report treated CDP20 as an external theorem that
immediately refuted the manuscript. That treatment is no longer
defensible after reading Section 10. The manuscript supplies a concrete
countermechanism, not a generic assertion of exception.

The appropriate present conclusion is narrower:

\begin{itemize}
\tightlist
\item
  The manuscript and the printed CDP20 conclusions are formally
  incompatible.
\item
  Section 10 gives a coherent explanation of why the printed proof does
  not cover the proposed nonnormal fibre.
\item
  This review found no internal defect in that explanation.
\item
  The explanation has not, by virtue of this review, been externally
  adjudicated. Acceptance would effectively publish a counterexample to
  CDP20 and should involve independent referees familiar with the
  original proof.
\end{itemize}

\section{Confirmed errors and
corrections}\label{confirmed-errors-and-corrections}

\subsection{Remark 7.28: quantifier
error}\label{remark-7.28-quantifier-error}

The manuscript's residue characterization is false as a biconditional
for a fixed triple. The triple

\[
(\ell_0,\ell_1,\ell_2)=(100,1,-1)
\]

has the listed residue pair, but

\[
p=12(100)-4(1)-3(-1)=1199\ne\pm1.
\]

The correct statement is:

\begin{quote}
For fixed admissible \(\ell_1,\ell_2\), there exists \(\ell_0\in\Z\)
such that \(|12\ell_0-4\ell_1-3\ell_2|=1\) if and only if the residue
pair is one of the two listed pairs.
\end{quote}

Alternatively, the remark should retain the equation involving
\(\ell_0\) and not reduce it to residues alone. This error is not used
in the chosen example and does not affect the main theorem.

\subsection{Proposition 4.7(iv): closed-image
justification}\label{proposition-4.7iv-closed-image-justification}

The statement that the section image is closed because it is ``proper
over the disc'' is not the correct proof. Replace it with the
closed-graph argument for a section into a Hausdorff total space or an
explicit quotient-chart proof. The conclusion remains valid.

\section{Major revision requirements}\label{major-revision-requirements}

The following revisions are necessary before the manuscript could
reasonably be accepted.

\subsection{1. Isolate and expand the two topology
engines}\label{isolate-and-expand-the-two-topology-engines}

Proposition 7.2 and Appendix B should each be independently readable.
The collapse proof needs a complete equivariance and continuity theorem,
preferably with a cell-level model. Appendix B needs a formal
nearby-cycle and specialization setup, precise integral references, and
a fully diagrammed saturation argument. The manuscript should state
exactly which conclusions are independent and which share inputs.

\subsection{2. Make analytic base change and duality
citation-complete}\label{make-analytic-base-change-and-duality-citation-complete}

Proposition 9.13 and Theorem 10.5 rely on the most delicate analytic
tools in the paper. For every use, state the morphism, sheaf, flatness
hypothesis, fibre dimension, and exact conclusion. In particular,
clarify top-degree base change at every scheme-theoretic fibre, duality
on the Gorenstein nonnormal surface \(W\), and the treatment of
nilpotent multiple fibres.

\subsection{3. Separate the CDP critique into a specialist-verifiable
note}\label{separate-the-cdp-critique-into-a-specialist-verifiable-note}

Section 10 should be presented with exact quotations, theorem numbers,
and a commutative diagram showing the failed reduction. The local
conductor construction should be stated as a general proposition
independent of the example. The authors should seek independent comments
from the authors of CDP20 or equivalent experts. This is not a demand
for deference to the literature; it is a demand proportionate to
overturning a published theorem.

\subsection{4. Add a complete convention and gluing
table}\label{add-a-complete-convention-and-gluing-table}

The signs, meridian orientations, deck generators, and translations
should be summarized in one table. The manuscript's sign lemma is
careful, but the information is dispersed. This is a preventable source
of catastrophic error.

\subsection{5. Strengthen provenance for a formal
submission}\label{strengthen-provenance-for-a-formal-submission}

Because the reviewed source is a reconstruction, a journal submission
must be tied to an author-approved version. The reconstruction audit can
remain supplementary evidence, but it cannot substitute for authorial
verification that every formula and correction reflects the intended
manuscript.

\subsection{6. Correct the two confirmed local
defects}\label{correct-the-two-confirmed-local-defects}

Remark 7.28 and the proof sentence in Proposition 4.7(iv) should be
corrected explicitly in a revision ledger.

\section{Computational audit}\label{computational-audit}

The accompanying script performs twenty-eight exact checks, all passing.
It uses integer matrices, Smith normal forms, determinantal divisors,
exterior powers, and symbolic polynomial identities. The checks include:

\begin{itemize}
\tightlist
\item
  \(T_1^3=T_2^4=I\), exact orders, \(T_1T_2T_0=I\), and \(N^2=0\);
\item
  the displayed contragredient matrices and their product relation;
\item
  the fixed lattices, the unimodular map \(B_0\), and
  invariance/signature of \(Q_0\);
\item
  the determinant and Smith normal form in Theorem 7.17;
\item
  all admissible congruence classes proving \(\gcd(p,12)=1\);
\item
  the relation-lattice primitivity in Lemma 7.19;
\item
  exterior-power monodromy, Wang ranks \((1,3,6,6,3,1)\), and the
  saturated image in degree two;
\item
  the rank, saturation, and primitive left kernel of the displayed
  \(\alpha_2\) matrix;
\item
  a concrete counterexample to the literal wording of Remark 7.28 and
  verification of the corrected existential version;
\item
  the expansion of \(\sigma_1\wedge\sigma_2\) and the formula for
  \(P_E\);
\item
  the \(g_0\) and \(g_2\) transformation identities in Lemma 9.3;
\item
  \(\eta^2=12\), the Gram matrix \(G_\eta\), and its determinant
  factorization;
\item
  explicit bookkeeping of the formal CDP20 conflict and the fact that
  Section 10, not arithmetic, is where applicability is disputed.
\end{itemize}

These checks validate formulas, not geometric existence. In particular,
a correct Smith normal form does not prove that the relation matrix is
the correct topological attaching matrix; that depends on the gluing and
sign arguments. The report maintains that distinction throughout.

\section{Final recommendation}\label{final-recommendation}

I recommend \textbf{major revision and resubmission}, with the following
editorial interpretation.

The manuscript should not be rejected merely by citing CDP20. The
complete text gives a concrete and technically credible challenge to the
proof step relevant to a nonnormal fibre. The earlier packet-only
rejection rationale is superseded.

The manuscript also should not be accepted solely because no flaw was
found in one extended review. It claims the first complex structure on
\(S^6\), uses a novel nonpolarized torus degeneration, depends on
delicate integral specialization, and would falsify printed conclusions
in the literature. The appropriate standard is independent replication
of the cusp topology and the Section 10 analytic mechanism.

Subject to that replication, the proof architecture is serious. It is
not a collection of unsupported claims: the construction is explicit,
the finite algebra is unusually detailed, two topology computations
agree, and the literature conflict is confronted directly. The confirmed
errors are local and repairable. The unresolved issues are high-risk
proof-verification issues, not demonstrated contradictions.

A revised manuscript should be sent to at least three specialist readers
covering: (1) toric and semi-abelian degenerations with nonpolarized
monodromy; (2) integral nearby cycles and topology of normal-crossing
degenerations; and (3) compact complex threefolds, analytic duality, and
the CDP98/CDP20 arguments. Acceptance should depend on affirmative
reports from those domains.

\section{References}\label{references}

Agricola, I., Bazzoni, G., Goertsches, O., Konstantis, P., \& Rollenske,
S. (2018). On the history of the Hopf problem. \emph{Differential
Geometry and its Applications, 57}, 1--9.

Ash, A., Mumford, D., Rapoport, M., \& Tai, Y.-S. (2010). \emph{Smooth
compactifications of locally symmetric varieties} (2nd ed.). Cambridge
University Press.

Banica, C., \& Stanasila, O. (1976). \emph{Algebraic methods in the
global theory of complex spaces}. Wiley.

Barth, W., Hulek, K., Peters, C., \& Van de Ven, A. (2004).
\emph{Compact complex surfaces} (2nd ed.). Springer.

Birkenhake, C., \& Lange, H. (2004). \emph{Complex abelian varieties}
(2nd ed.). Springer.

Bredon, G. E. (1997). \emph{Sheaf theory} (2nd ed.). Springer.

Campana, F., Demailly, J.-P., \& Peternell, T. (1998). The algebraic
dimension of compact complex threefolds with vanishing second Betti
number. \emph{Compositio Mathematica, 112}, 77--91.

Campana, F., Demailly, J.-P., \& Peternell, T. (2020). Corrigendum: The
algebraic dimension of compact complex threefolds with vanishing second
Betti number. \emph{Compositio Mathematica, 156}, 679--696.

Clemens, C. H. (1977). Degeneration of Kahler manifolds. \emph{Duke
Mathematical Journal, 44}, 215--290.

Deligne, P. (1973). Le formalisme des cycles evanescents. In P. Deligne
\& N. Katz (Eds.), \emph{Groupes de monodromie en geometrie algebrique
(SGA 7 II)} (Lecture Notes in Mathematics 340, pp.~82--115). Springer.

Dimca, A. (2004). \emph{Sheaves in topology}. Springer.

Fulton, W. (1993). \emph{Introduction to toric varieties}. Princeton
University Press.

Hatcher, A. (2002). \emph{Algebraic topology}. Cambridge University
Press.

Kashiwara, M., \& Schapira, P. (1990). \emph{Sheaves on manifolds}.
Springer.

Kervaire, M. A., \& Milnor, J. W. (1963). Groups of homotopy spheres I.
\emph{Annals of Mathematics, 77}, 504--537.

Lange, H., \& Birkenhake, C. (1999). \emph{Complex tori}. Birkhauser.

Mumford, D. (1972). An analytic construction of degenerating abelian
varieties over complete rings. \emph{Compositio Mathematica, 24},
239--272.

Ramis, J.-P., \& Ruget, G. (1970). Complexe dualisant et theoremes de
dualite en geometrie analytique complexe. \emph{Publications
Mathematiques de l'IHES, 38}, 77--91.

Ramis, J.-P., Ruget, G., \& Verdier, J.-L. (1971). Dualite relative en
geometrie analytique complexe. \emph{Inventiones Mathematicae, 13},
261--283.

Smale, S. (1961). Generalized Poincare's conjecture in dimensions
greater than four. \emph{Annals of Mathematics, 74}, 391--406.

\section{Appendix A: audit disposition by central
claim}\label{appendix-a-audit-disposition-by-central-claim}

\begin{longtable}[]{@{}
  >{\raggedright\arraybackslash}p{(\columnwidth - 6\tabcolsep) * \real{0.2500}}
  >{\raggedright\arraybackslash}p{(\columnwidth - 6\tabcolsep) * \real{0.2500}}
  >{\raggedright\arraybackslash}p{(\columnwidth - 6\tabcolsep) * \real{0.2500}}
  >{\raggedright\arraybackslash}p{(\columnwidth - 6\tabcolsep) * \real{0.2500}}@{}}
\toprule\noalign{}
\begin{minipage}[b]{\linewidth}\raggedright
Central claim
\end{minipage} & \begin{minipage}[b]{\linewidth}\raggedright
Audit disposition
\end{minipage} & \begin{minipage}[b]{\linewidth}\raggedright
Principal dependency
\end{minipage} & \begin{minipage}[b]{\linewidth}\raggedright
Required follow-up
\end{minipage} \\
\midrule\noalign{}
\endhead
\bottomrule\noalign{}
\endlastfoot
Smooth compact complex threefold \(X\) exists & No fatal defect found &
Sections 2--6 & Independent review of infinite toric quotient and
gluing \\
\(\pi_1(X)=1\) & Exact algebra verified, conditional on attaching
relations & Section 7.5 & Orientation/gluing table \\
\(H^*(X;\Z)=H^*(S^6;\Z)\) & Two routes agree; not fully certified &
Sections 7.2--7.7, Appendix B & Replicate collapse and nearby-cycle
specialization \\
\(X\cong_{\mathrm{diff}}S^6\) & Standard conditional consequence &
Section 8 & None beyond topology verification \\
\(a(X)=1\) & Coherent; period algebra verified & Section 9.1 & Expand
torus algebraic-dimension criterion \\
Direct images and \(K_X\) & Plausible & Sections 9.2--9.4 & Exact
base-change/duality statements \\
\(c_3(X)=2\) & Conditional on topology & Sections 7--9 & No separate
arithmetic issue \\
CDP20 does not apply & Substantive rebuttal; not externally adjudicated
& Section 10 & Specialist independent review and literature response \\
\end{longtable}

\section{Appendix B: supersession
notice}\label{appendix-b-supersession-notice}

This report supersedes the prior packet-06 review. That earlier review
correctly identified the formal incompatibility between the claims
visible on pages 56--66 and the printed CDP20 conclusions, but it did
not have the complete manuscript. In particular, it lacked Section 10,
Appendix B, and the complete construction. Its recommendation of
rejection on the ground of an unrebutted CDP20 contradiction is
withdrawn. The exact packet-06 arithmetic findings remain useful and
have been incorporated into the present full-manuscript checks.

\section{Appendix C: reproducibility}\label{appendix-c-reproducibility}

The review bundle contains:

\begin{itemize}
\tightlist
\item
  \texttt{report/full\_mathematical\_peer\_review.pdf}, the primary
  report;
\item
  \texttt{report/full\_mathematical\_peer\_review.md} and \texttt{.tex},
  editable sources;
\item
  \texttt{audit/issue\_ledger.csv} and \texttt{.json};
\item
  \texttt{audit/theorem\_dependency\_matrix.csv};
\item
  \texttt{audit/exact\_checks.py}, \texttt{exact\_checks\_results.json},
  and execution logs;
\item
  \texttt{evidence/archive\_verify\_stdout.txt} and reconstruction-audit
  extracts;
\item
  \texttt{evidence/source\_page\_line\_map.csv} and \texttt{.md};
\item
  \texttt{MANIFEST.sha256} and a review receipt.
\end{itemize}

To rerun the arithmetic audit from the bundle root:

\begin{Shaded}
\begin{Highlighting}[]
\ExtensionTok{python}\NormalTok{ audit/exact\_checks.py}
\end{Highlighting}
\end{Shaded}

To rebuild the report, use Pandoc with XeLaTeX or compile the included
generated LaTeX source. The PDF was rendered at 200 dpi and preflighted
before packaging.

\end{document}
